% META: {"id": "1NSI-PM-ME-002", "chapitre": "1NSI-PROJET-METHODES", "type_objet": "methode", "capacites": ["BO-PREAMBULE-DEMARCHE-DE-PROJET"], "capacites_codes": ["C2"], "methodes": ["M2"], "mode_creation": "ex_nihilo", "sources_inspiration": [], "status": "needs_review"}
% BEGIN-VERIFY
% def avancement(jalons):
%     if len(jalons) == 0:
%         raise ValueError("au moins un jalon est requis")
%     nb_termines = sum(1 for j in jalons if j["termine"])
%     return nb_termines / len(jalons) * 100
% def prochain_jalon(jalons):
%     for j in jalons:
%         if not j["termine"]:
%             return j["nom"]
%     return None
% jalons = [
%     {"nom": "Cahier des charges ecrit", "critere": "liste des fonctionnalites relue par le professeur", "termine": True},
%     {"nom": "Deplacement teste", "critere": "le personnage bouge avec les fleches, 3 tests passent", "termine": True},
%     {"nom": "Score affiche", "critere": "le score augmente a chaque piece ramassee", "termine": False},
%     {"nom": "Presentation repetee", "critere": "demonstration de 2 minutes chronometree", "termine": False},
% ]
% assert avancement(jalons) == 50.0
% assert prochain_jalon(jalons) == "Score affiche"
% assert all(len(j["critere"]) > 0 for j in jalons)
% END-VERIFY
\begin{fichemethode}{M2}{De l'idée aux jalons : un cahier des charges court, des étapes vérifiables, un suivi}

\textbf{Quand l'utiliser ?} Au démarrage d'un projet en équipe, ou quand l'énoncé demande de
découper un projet, d'en estimer l'avancement, ou de juger si une décision d'équipe est
raisonnable.

\textbf{Pas à pas :}
\begin{enumerate}
  \item \textbf{Écris le cahier des charges en une liste de fonctionnalités vérifiables},
        une par ligne : « le personnage se déplace avec les flèches », pas « le jeu est
        bien ». Chaque ligne doit pouvoir être déclarée vraie ou fausse en la testant.
  \item \textbf{Sépare l'essentiel de l'optionnel.} L'ambition doit rester raisonnable pour
        une équipe de $2$ à $4$ élèves : un projet modeste, complet et testé vaut mieux qu'un
        projet ambitieux inachevé. Les options viennent après, si le temps le permet.
  \item \textbf{Découpe en $4$ à $6$ jalons ordonnés}, chacun avec un \emph{critère de
        validation} exécutable (un test qui passe, une démonstration). Le premier jalon est
        le cahier des charges lui-même ; le dernier est la présentation.
  \item \textbf{Répartis les tâches} par fonctionnalité ou par couche (données, algorithmes,
        interface), et fixe des points d'étape avec le professeur.
  \item \textbf{Suis l'avancement} avec la structure \lstinline{jalons} du cours :
        \lstinline{avancement} donne un pourcentage, \lstinline{prochain_jalon} dit où
        concentrer l'effort. Lis ce pourcentage avec prudence : il compte les jalons, pas
        leur charge de travail.
\end{enumerate}

\textbf{Exemple rédigé.} Un petit jeu de plateforme. Cahier des charges essentiel : le
personnage se déplace ; il ramasse des pièces ; le score s'affiche. Optionnel : plusieurs
niveaux, une musique. Jalons : cahier des charges écrit (critère : liste relue par le
professeur) ; déplacement testé (trois tests passent) ; score affiché (le score augmente à
chaque pièce) ; présentation répétée (démonstration de $2$ minutes). Deux jalons terminés
sur quatre : \lstinline{avancement(jalons)} vaut $50{,}0$ et \lstinline{prochain_jalon}
renvoie \lstinline{"Score affiche"}. Si les niveaux multiples ne sont pas faits le jour de
la présentation, le projet est quand même complet : ils étaient optionnels.

\textbf{Erreur fréquente.} Des jalons invérifiables (« tout coder », « que ça marche »), qui
ne disent jamais où l'on en est. Ou croire qu'un jalon garantit l'absence de bug : seul un
critère de validation exécuté peut le dire.

\end{fichemethode}
